%------------------------------------------------------------------------------%

\usepackage{integracao_e_series}

\title{Séries de Potências -- Convergência}

%------------------------------------------------------------------------------%
\begin{document}

\addtitle

%------------------------------------------------------------------------------%
\section{Convergência Ponto a Ponto}

%------------------------------------------------------------------------------%
\begin{frame}{Convergência das Séries de Potências}
  \begin{itemize}[<+->]
    \item \emph{Convergência Trivial:} Toda Série de Potências converge em seu centro\\[5mm]
    \item Tipicamente empregamos o Teste da \emph{Razão} ou da \emph{Raiz}
  \end{itemize}
\end{frame}

%------------------------------------------------------------------------------%
\addtocounter{exemplo}{1}
\begin{frame}{Exemplo \theexemplo}
  Considerando a série
  \[
    \sum_{n=1}^{\infty} \frac{x^n}{n}
    = x + \frac{x^2}{2}
        + \frac{x^3}{3}
        + \frac{x^4}{4}
        + \frac{x^5}{5}
        + \cdots
  \]
  \pause
  Aplicando o teste da razão para $x\neq 0$
  \[
    \rho   = \lim_{n\to\infty} \frac{\;\abs{u_{n+1}}\;}{\abs{u_n}}
    \pause = \lim_{n\to\infty} \frac{\abs{x^{n+1}}}{\;n+1\;} \frac{n}{\;\abs{x^n}\;}
    \pause = \lim_{n\to\infty} \frac{n}{\;n+1\;} \abs{x}
    \pause = \abs{x} \lim_{n\to\infty} \frac{n}{\;n+1\;}
    \pause = \abs{x}
  \]
  \pause
  \emph{Converge absolutamente} para todo $x$ no intervalo aberto $(-1,\,1\,)$
\end{frame}

%------------------------------------------------------------------------------%
\addtocounter{exemplo}{1}
\begin{frame}{Exemplo \theexemplo}
  Considerando a série
  \[
    \sum_{n=0}^{\infty} \frac{x^n}{n!}
    = 1 + x
        + \frac{x^2}{2}
        + \frac{x^3}{3!}
        + \frac{x^4}{4!}
        + \frac{x^5}{5!}
        + \cdots
  \]
  \pause
  Aplicando o teste da razão para $x\neq 0$
  \[
    \rho   = \lim_{n\to\infty} \frac{\;\abs{u_{n+1}}\;}{\abs{u_n}}
    \pause = \lim_{n\to\infty} \frac{\abs{x^{n+1}}}{\;(n+1)!\;} \frac{n!}{\;\abs{x^n}\;}
    \pause = \lim_{n\to\infty} \frac{1}{\;n+1\;} \abs{x}
    \pause = \abs{x} \lim_{n\to\infty} \frac{1}{\;n+1\;}
    \pause = 0
  \]
  \pause
  \emph{Converge absolutamente} para todo $x$ no intervalo $(-\infty,\,\infty\,)$
\end{frame}

%------------------------------------------------------------------------------%
\section{Convergência de uma Série de Potências}

%------------------------------------------------------------------------------%
\begin{frame}{Convergência de uma Série de Potências -- 1}
  Dada a série de potências
  \[
    \sum_{n=0}^\infty c_n(x-a)^n
  \]

  \pause
  \vspace{\baselineskip}
  Se ela converge em \({\color{blue}L}=x-a\neq 0\) \\[1mm] \pause
  então ela converge absolutamente para todo $x$ com
  \(
    \abs{x-a} < \abs{{\color{blue}L}}
  \)

  \pause
  \vspace{\baselineskip}
  Se ela diverge em \({\color{blue}M}=x-a\) \\[1mm] \pause
  então ela diverge para todo $x$ tal que
  \(
    \abs{x-a}>\abs{{\color{blue}M}}
  \)

\end{frame}

%------------------------------------------------------------------------------%
\addtocounter{exemplo}{1}
\begin{frame}{Exemplo \theexemplo}
  Determine para quais valores de $x$ a série
  \[
    \sum_{n=0}^\infty (-1)^n\frac{(x-2)^n}{n}
  \]
  é convergente
\end{frame}

%------------------------------------------------------------------------------%
\begin{frame}{Exemplo \theexemplo}
  Quando \(x-2=1\), isso é, \(x=3\)
  \pause
  \[
    \sum_{n=0}^\infty (-1)^n\frac{(x-2)^n}{n}
    \;=\;
    \sum_{n=0}^\infty \frac{(-1)^n}{n}
  \]

  \pause
  Série Harmônica Alternada, que converge condicionalmente

  \pause
  \vspace{\baselineskip}
  A série converge para
  \( L = 1 \)\pause,
  então, \emph{converge absolutamente} para todo \(x\) tal que
  \[
    \abs{x-2} < \abs{L} = 1
  \]
  \pause
  ou seja, no intervalo
  \(
    (1, 3)
  \)
\end{frame}

%------------------------------------------------------------------------------%
\begin{frame}{Exemplo \theexemplo}
  Quando \(x-2=-1\), isso é, \(x=1\)
  \pause
  \[
    \sum_{n=0}^\infty (-1)^n\frac{(x-2)^n}{n}
    \pause
    \;=\;
    \sum_{n=0}^\infty \frac{(-1)^n(-1)^n}{n}
    \pause
    \;=\;
    \sum_{n=0}^\infty \frac{1}{n}
  \]

  \pause
  Série Harmônica, que diverge

  \pause
  \vspace{\baselineskip}
  A série diverge para
  \( M = -1 \)\pause,
  então diverge para todo \(x\) tal que
  \[
    \abs{x-2} > \abs{M} = 1
  \]
  \pause
  ou seja, em
  \(
    (-\infty, 1) \cup (3, \infty)
  \)
\end{frame}

%------------------------------------------------------------------------------%
\begin{frame}{Exemplo \theexemplo}
  A série converge absolutamente no intervalo
  \[
    (1, 3)
  \]

  \pause
  A série converge no intervalo
  \[
    (1, 3\,]
  \]
\end{frame}

%------------------------------------------------------------------------------%
\begin{frame}{Convergência de uma Série de Potências -- 2}
  Sobre a convergência de
  \(\quad\sum_{n=0}^\infty c_n(x-a)^n\quad\)

  \pause
  \vspace{\baselineskip}
  \begin{enumerate}[<+->]
    \item Existe um número ${\color{blue}R>0}$ tal que a série \\[1mm]
          \tabto{10mm} converge absolutamente para \tabto{70mm} \(\abs{x-a}<{\color{blue}R}\) \\
          \tabto{10mm} diverge para                \tabto{70mm} \(\abs{x-a}>{\color{blue}R}\)
    \vspace{\baselineskip}
    \item A série converge para todo $x$, nesse caso escrevemos ${\color{blue}R=\infty}$
    \vspace{\baselineskip}
    \item A série diverge para todo $x\neq a$, nesse caso ${\color{blue}R=0}$
  \end{enumerate}
\end{frame}

%------------------------------------------------------------------------------%
\begin{frame}{Raio e Intervalo de Convergência}
  \begin{itemize}[<+->]
    \item \({\color{blue}R}\) é o de \emph{raio de convergência} \\[\baselineskip]
    \item o \emph{intervalo de convergência} é o intervalo centrado em \(a\) com raio ${\color{blue}R}$
          \[
            I = \big( a-R, a+R \big)
          \]
  \end{itemize}

  \action<+->{
  Os extremos do intervalo de convergência, $x=a\pm {\color{blue}R}$, \\
  podem ser abertos ou fechados}
\end{frame}

%------------------------------------------------------------------------------%
\begin{frame}{Testando a Convergência}
  \begin{enumerate}[<+->]
  \item Utilize o teste da razão (ou da raiz) para determinar o raio de convergência
  \vspace{\baselineskip}
  \item Se \(R\) for finito teste os extremos do intervalo
        \begin{itemize}
          \item teste da comparação
          \item teste da integral
          \item teste da série alternada
        \end{itemize}
  \end{enumerate}
\end{frame}

%------------------------------------------------------------------------------%
\addtocounter{exemplo}{1}
\begin{frame}{Exemplo \theexemplo}
  Analise a convergência da série de potências
  \[
    \sum_{n=0}^\infty \frac{(x-2)^n}{10^n}
  \]
\end{frame}

%------------------------------------------------------------------------------%
\begin{frame}{Exemplo \theexemplo}
  Termo geral da série
  \(
    a_n = \frac{(x-2)^n}{10^n}
  \)

  \pause
  \vspace{\baselineskip}
  Aplicando o Teste da Raiz no módulo
  \pause
  \[
    \sqrt[n]{\abs{a_n}}                    \pause =
    \sqrt[n]{\abs{\frac{(x-2)^n}{10^n}}}   \pause =
    \sqrt[n]{\abs{\frac{x-2}{10}}^n}       \pause =
    \frac{\abs{x-2}}{10}
  \]
  \pause
  Calculando o limite
  \[
     \rho
     = \lim_{n\to\infty} \frac{\abs{x-2}}{10}
     \pause
     = \frac{\abs{x-2}}{10}
  \]
\end{frame}

%------------------------------------------------------------------------------%
\begin{frame}{Exemplo \theexemplo}
\begin{columns}
\column{0.3\linewidth}
  \onslide<+->{A série converge quando}
  \begin{align*}
    \onslide<1->{\rho                 & < 1  \\[2mm]}
    \onslide<+->{\frac{\abs{x-2}}{10} & < 1  \\[2mm]}
    \onslide<+->{\abs{x-2}            & < 10}
  \end{align*}
\column{0.5\linewidth}
  \onslide<+->{\emph{Centro do Intervalo de Convergência}
  \[
    x = 2
  \]}
  \onslide<+->{\emph{Raio de Convergência}
  \[
    R = 10
  \]}
  \onslide<+->{\emph{Interior do Intervalo de Convergência}
  \[
    (-8, 12)
  \]}
  \onslide<+->{Falta analisar os extremos}
\end{columns}
\end{frame}

%------------------------------------------------------------------------------%
\begin{frame}{Exemplo \theexemplo}
  Considerando \(x = -8\), isto é, \(x-2=-10\)
  \pause
  \[
    \sum_{n=0}^\infty \frac{(x-2)^n}{10^n} \pause \;=\;
    \sum_{n=0}^\infty \frac{(-10)^n}{10^n} \pause \;=\;
    \sum_{n=0}^\infty (-1)^n
  \]
  \pause
  Diverge pelo Teste da Divergência
\end{frame}

%------------------------------------------------------------------------------%
\begin{frame}{Exemplo \theexemplo}
  Considerando \(x = 12\), isto é, \(x-2=10\)
  \pause
  \[
    \sum_{n=0}^\infty \frac{(x-2)^n}{10^n} \pause \;=\;
    \sum_{n=0}^\infty \frac{10^n}{10^n}    \pause \;=\;
    \sum_{n=0}^\infty 1
  \]
  \pause
  Diverge pelo Teste da Divergência
\end{frame}

%------------------------------------------------------------------------------%
\begin{frame}{Exemplo \theexemplo}
  \onslide<+->{\emph{Centro do Intervalo de Convergência}
  \[
    x = 2
  \]}
  \onslide<+->{\emph{Raio de Convergência}
  \[
    R = 10
  \]}
  \onslide<+->{\emph{Intervalo de Convergência}
  \[
    (-8, 12)
  \]}
\end{frame}

%------------------------------------------------------------------------------%
\addtocounter{exemplo}{1}
\begin{frame}{Exemplo \theexemplo}
  Analise a convergência da série de potências
  \[
    \sum_{n=0}^\infty \frac{(-x)^n}{n!}
  \]
\end{frame}

%------------------------------------------------------------------------------%
\begin{frame}{Exemplo \theexemplo}
  Teste da Razão para \(x \neq 0\)
  \pause
  \[
    \abs{\frac{a_{n+1}}{a_n}}                                \pause =
    \abs{\frac{(-x)^{n+1}}{(n+1)!}\,\frac{n!}{(-x)^n}}       \pause =
    \frac{\abs{(-x)^{n+1}}}{(n+1)!}\,\frac{n!}{\abs{(-x)}^n} \pause =
    \frac{\abs{x}^{n+1}}{(n+1)n!}\,\frac{n!}{\abs{x}^n}      \pause =
    \frac{\abs{x}}{n+1}
  \]

  \pause
  Calculando o limite
  \[
    \rho
    = \lim_{n\to\infty} \frac{\abs{x}}{n+1}
    \pause
    = \abs{x} \lim_{n\to\infty} \frac{1}{n+1}
    \pause
    = 0
  \]

  \pause
  Portanto a série converge absolutamente para todo \(x\in\R\)

  \pause
  \vspace{0.5\baselineskip}
  Raio de Convergência \(R = \infty\)
\end{frame}

%------------------------------------------------------------------------------%
\addtocounter{exemplo}{1}
\begin{frame}{Exemplo \theexemplo}
  Analise a convergência da série de potências
  \[
    \sum_{n=0}^\infty \abs{n^n x^n}
  \]
\end{frame}

%------------------------------------------------------------------------------%
\begin{frame}{Exemplo \theexemplo}
  Teste da Raiz para \(x\neq 0\)
  \pause
  \[
    \sqrt[n]{\abs{n^nx^n}} =
    n\abs{x}
  \]
  \pause
  \vspace{-\baselineskip}
  \[
    \rho = \lim_{n\to\infty} \sqrt[n]{\abs{n^nx^n}} = \infty
  \]

  \pause
  \vspace{0.7\baselineskip}
  A série diverge para todo \(x\neq 0\)

  \pause
  \vspace{0.5\baselineskip}
  Raio de Convergência é $R=0$
\end{frame}

%------------------------------------------------------------------------------%
\begin{frame}{\contentsname}
  \tableofcontents
\end{frame}

%------------------------------------------------------------------------------%
\end{document}
%------------------------------------------------------------------------------%
